\documentclass[11pt,a4paper]{ctexart}
\usepackage[margin=2.2cm]{geometry}
\usepackage{amsmath,amssymb,booktabs,graphicx,float,xcolor,adjustbox,hyperref}
\usepackage[section]{placeins}
\hypersetup{colorlinks=true,linkcolor=blue!50!black,citecolor=blue!50!black,urlcolor=blue!50!black}
\title{各向异性张量系数：BoomerAMG 变差的例子，以及 FD-LGF 的表现}
\author{温州肯恩大学 \quad Nonlinear Poisson 项目（H34）}
\date{2026 年 11 月 3 日}

\begin{document}
\maketitle

\begin{abstract}
H33 表明：标量系数、通量差分得到的 M-矩阵上，BoomerAMG 极其稳健。本报告改用张量系数
$-\nabla\cdot(K\nabla u)+qu=f$，$K=\chi\,ee^{T}+(I-ee^{T})$（各向异性比 $\chi$），
标准对称 9 点格式离散。只要 $K_{12}\neq0$，矩阵就有正的非对角元，不再是 M-矩阵。
\textbf{结论：}(1) 这里确实找到了 BoomerAMG 变差的例子：圆形磁力线（circ）、局部旋转各向异性圆盘（disk）和斜断层带（band）上，
它的迭代数随网格加密持续增长，$n=1023$、$\chi\ge10^4$ 时分别要 144--241、98--218、56--65 次；常数旋转只是轻微增长（$\le33$ 次）。
(2) 各向异性局限在小区域（disk、band）时，FD-LGF（各向同性背景 + 各向异性区域上的精确 LU）只要 4--8 次迭代，并且与 $n$ 无关；
时间在 band 上是 BoomerAMG 的 0.29--0.31 倍，在 disk 上是 0.53--0.66 倍。
(3) 各向异性遍布全域（rot30、rot45、circ）时，FD-LGF 只有网格对齐的背景可用，条件数 $\sim\chi$，要 95--1000 次迭代，比 BoomerAMG 慢 1.8--6 倍。
(4) GMG 的迭代数是 BoomerAMG 的 2--5 倍以上。
预注册的 5 条预测中，1、2、4、5 成立；第 3 条在 band 上成立，在 disk 上时间未达到预测的 0.5 倍。
\end{abstract}

\section{问题与离散}
能量形式
\[
E(u)=\sum_{x\text{-面}}K_{11}^{f}(\delta_xu)^2+\sum_{y\text{-面}}K_{22}^{f}(\delta_yu)^2+2\sum_{\text{节点}}K_{12}(c_xu)(c_yu)+\sum qu^2 ,
\]
面上的值取相邻两节点的算术平均，$c_x,c_y$ 为中心差分；$A=\tfrac12\nabla^2E$ 就是标准对称 9 点格式。
$K=aI$ 时它与原来的 5 点通量格式逐元一致（已测试）。
由 $(c_xu)^2\le$ 相邻两个 $(\delta_xu)^2$ 的平均和 $K_{12}^2\le K_{11}K_{22}$，用 Cauchy--Schwarz 可得 $E(u)\ge\sum qu^2$，所以 $A$ 对称正定。
制造解与 H33 相同：$u^*=\sin\pi x\sin\pi y$，$f=A_hu^*$，$q=1+0.2\cos2\pi x\cos2\pi y$。

\begin{figure}[H]\centering
\includegraphics[width=\textwidth]{gallery.png}
\caption{六类系数，蓝色短线是强传导方向 $e$，灰色是各向异性区域（$\chi=10^4$ 示意）。
aligned 为网格对齐的对照；rot30、rot45 为常数旋转；circ 为绕中心的圆形磁力线（磁化等离子体型，arXiv:2301.13351）；
disk 为各向同性背景中半径 0.2 的 30° 圆盘；band 为宽 0.04 的 30° 斜断层带，沿带传导强。}
\end{figure}

\paragraph{方法。}BoomerAMG（PETSc 默认 $\theta=0.25$，以及 $\theta=0.5$）、GMG（$\omega=4/5$）、FD-LGF。
FD-LGF 的背景取对角张量 $\mathrm{diag}(k_1(x),k_2(x))$，用 DST-$y$ 加三对角-$x$ 精确求逆：
aligned 的背景就是 $K$ 本身；rot30、rot45、circ 取 $\mathrm{diag}(\overline{K_{11}},\overline{K_{22}})$，没有局部部分；
disk、band 取各向同性背景 $I$，在各向异性区域外扩宽度为 $w$ 的环后做精确 LU。
PCG 停止规则 $\|r\|/\|f\|\le10^{-10}$，单线程计时。停止后继续迭代（不计时、不计入迭代数）只记录误差，用来定义 $k_8$。

\section{确认性结果（$n=1023$，约 $10^6$ 未知量）}
\begin{table}[H]\centering\small
\caption{$n=1023$：PCG 迭代次数（各 $\chi$ 的范围）}\label{tab:its}
\begin{tabular}{lccccc}
\toprule
系数 & BAMG $\theta{=}0.25$ & BAMG $\theta{=}0.5$ & GMG & FD-LGF $w{=}\frac1{16}$ & FD-LGF $w{=}\frac1{64}$ \\
\midrule
aligned 对齐（对照） & 3--8 & 3--8 & 58--1000 & 2--3 & -- \\
rot30 旋转 30° & 12--28 & 10--19 & 43--111 & 95--244 & -- \\
rot45 旋转 45° & 13--33 & 13--25 & 40--104 & 112--287 & -- \\
circ 圆形磁力线 & 23--191 & 29--241 & 72--718 & 131--1000 & -- \\
disk 局部旋转圆盘 & 20--136 & 28--218 & 50--732 & 4--7 & 7--11 \\
band 斜断层带 & 22--59 & 24--65 & 39--167 & 6--8 & 11--14 \\
\bottomrule
\end{tabular}
\end{table}

\begin{table}[H]\centering\small
\caption{$n=1023$：误差降到 $10^{-8}$ 所需迭代数}\label{tab:k8}
\begin{tabular}{lccccc}
\toprule
系数 & BAMG $\theta{=}0.25$ & BAMG $\theta{=}0.5$ & GMG & FD-LGF $w{=}\frac1{16}$ & FD-LGF $w{=}\frac1{64}$ \\
\midrule
aligned 对齐（对照） & 2--5 & 2--5 & 28--407 & 1--2 & -- \\
rot30 旋转 30° & 7--17 & 6--12 & 23--56 & 50--125 & -- \\
rot45 旋转 45° & 8--20 & 8--16 & 21--52 & 61--156 & -- \\
circ 圆形磁力线 & 14--171 & 19--217 & 39--646 & 82--805 & -- \\
disk 局部旋转圆盘 & 14--129 & 19--211 & 30--682 & 6--7 & 7--10 \\
band 斜断层带 & 14--43 & 16--49 & 23--131 & 7--8 & 11--16 \\
\bottomrule
\end{tabular}
\end{table}

\begin{table}[H]\centering\small
\caption{$n=1023$：总时间 setup+solve（s，单线程）}\label{tab:time}
\begin{tabular}{lccccc}
\toprule
系数 & BAMG $\theta{=}0.25$ & BAMG $\theta{=}0.5$ & GMG & FD-LGF $w{=}\frac1{16}$ & FD-LGF $w{=}\frac1{64}$ \\
\midrule
aligned 对齐（对照） & 0.67--1.18 & 0.67--1.22 & 3.70--52.12 & 0.09--0.15 & -- \\
rot30 旋转 30° & 2.52--5.35 & 1.66--2.52 & 3.49--7.63 & 6.39--20.23 & -- \\
rot45 旋转 45° & 1.91--4.22 & 3.00--4.80 & 3.25--7.47 & 5.91--13.64 & -- \\
circ 圆形磁力线 & 4.35--26.96 & 5.23--34.20 & 5.39--51.16 & 6.45--54.61 & -- \\
disk 局部旋转圆盘 & 2.40--11.00 & 3.52--21.02 & 3.02--45.33 & 5.39--5.88 & 3.91--4.04 \\
band 斜断层带 & 2.31--4.82 & 2.84--6.25 & 2.18--9.55 & 1.38--1.53 & 1.01--1.19 \\
\bottomrule
\end{tabular}
\end{table}

\begin{table}[H]\centering\small
\caption{$n=1023$：setup 时间（s）}\label{tab:setup}
\begin{tabular}{lccccc}
\toprule
系数 & BAMG $\theta{=}0.25$ & BAMG $\theta{=}0.5$ & GMG & FD-LGF $w{=}\frac1{16}$ & FD-LGF $w{=}\frac1{64}$ \\
\midrule
aligned 对齐（对照） & 0.45--0.51 & 0.45--0.57 & 0.27--0.32 & 0.02 & -- \\
rot30 旋转 30° & 0.89--0.92 & 0.69--0.75 & 0.29--0.34 & 0.02 & -- \\
rot45 旋转 45° & 0.78--0.86 & 0.92--0.96 & 0.32--0.42 & 0.02 & -- \\
circ 圆形磁力线 & 1.10--1.21 & 1.03--1.26 & 0.33--0.35 & 0.02 & -- \\
disk 局部旋转圆盘 & 0.68--0.76 & 0.64--0.74 & 0.25--0.27 & 4.61--5.24 & 2.86--3.02 \\
band 斜断层带 & 0.65--0.66 & 0.64--0.65 & 0.23--0.26 & 0.87--0.88 & 0.32--0.33 \\
\bottomrule
\end{tabular}
\end{table}

\begin{table}[H]\centering\small
\caption{$n=1023$：停止时的代数误差（各 $\chi$ 最大值）}\label{tab:err}
\begin{tabular}{lccccc}
\toprule
系数 & BAMG $\theta{=}0.25$ & BAMG $\theta{=}0.5$ & GMG & FD-LGF $w{=}\frac1{16}$ & FD-LGF $w{=}\frac1{64}$ \\
\midrule
aligned 对齐（对照） & $9.5{\times}10^{-13}$ & $9.5{\times}10^{-13}$ & $7.0{\times}10^{-12}$ & $4.1{\times}10^{-13}$ & -- \\
rot30 旋转 30° & $3.1{\times}10^{-13}$ & $7.0{\times}10^{-13}$ & $1.1{\times}10^{-12}$ & $7.5{\times}10^{-13}$ & -- \\
rot45 旋转 45° & $4.9{\times}10^{-13}$ & $1.5{\times}10^{-12}$ & $9.2{\times}10^{-13}$ & $7.2{\times}10^{-13}$ & -- \\
circ 圆形磁力线 & $2.1{\times}10^{-9}$ & $1.5{\times}10^{-9}$ & $1.3{\times}10^{-9}$ & $4.2{\times}10^{-3}$ & -- \\
disk 局部旋转圆盘 & $7.7{\times}10^{-9}$ & $3.5{\times}10^{-9}$ & $2.7{\times}10^{-9}$ & $1.8{\times}10^{-6}$ & $1.4{\times}10^{-7}$ \\
band 斜断层带 & $4.6{\times}10^{-12}$ & $6.7{\times}10^{-12}$ & $4.0{\times}10^{-11}$ & $1.6{\times}10^{-7}$ & $2.4{\times}10^{-7}$ \\
\bottomrule
\end{tabular}
\end{table}

\begin{figure}[H]\centering
\includegraphics[width=\textwidth]{scaling_1e+04.png}
\caption{迭代数随 $n$ 的变化（$\chi=10^4$）。BoomerAMG 在 circ、disk、band 上随 $n$ 增长；FD-LGF 在 disk、band 上保持平坦。}
\end{figure}
\begin{figure}[H]\centering
\includegraphics[width=\textwidth]{scaling_1e+06.png}
\caption{同上，$\chi=10^6$。}
\end{figure}

\section{预注册预测的检验}
方案在运行前锁定（哈希 bf86f18d，2026-11-03T13:50:23Z）。$n=511$ 的参考值来自用同一份锁定代码重跑的扫描，其迭代数与第一次扫描逐行相同。
\begin{enumerate}
\item \textbf{成立。}BoomerAMG 的迭代数从 $n=511$ 到 $n=1023$：circ 在 $\theta=0.25$ 时为 $96\to144$、$108\to191$，
在 $\theta=0.5$ 时为 $113\to192$、$128\to241$，都满足 $\ge1.3$ 倍；disk 为 $74\to98$、$85\to136$ 和 $101\to175$、$113\to218$；
band 为 $47\to56$、$48\to59$ 和 $45\to62$、$46\to65$。circ 在 $\theta=0.25$ 时 $>120$ 次。
\item \textbf{成立。}aligned 上 FD-LGF 只要 2--3 次，时间是 BoomerAMG 的 0.12--0.13 倍；
rot30、rot45、circ 上 FD-LGF 的时间是最快 BoomerAMG 的 1.8--8.0 倍（预测 $\ge1.5$ 倍）。
\item \textbf{部分成立。}disk、band 上 FD-LGF（$w=1/16$）的 $k_8=6$--8，不超过 $n=511$ 时 +1，这部分成立；
时间在 band 上是 0.29--0.31 倍（成立），在 disk 上是 0.53--0.66 倍，没有达到预测的 $\le0.5$。
disk 的 LU 区域占网格的 20\%，setup 占了大部分时间；$w=1/64$ 时是 0.36--0.48 倍。
\item \textbf{成立。}GMG 的迭代数在 aligned 上是 BoomerAMG 的 54--333 倍，在 circ 上是 3.0--3.8 倍，在 disk 上是 4.0--5.4 倍。
\item \textbf{成立。}$\chi=10^6$ 时 FD-LGF 按残差规则停止，此时代数误差在 disk 上为 $2\times10^{-6}$，在 band 上为 $2\times10^{-7}$；
再迭代 3 次和 2 次就降到 $10^{-8}$。
\end{enumerate}

\paragraph{计时说明。}这次确认性运行期间，有其他实验在同时运行（各行记录的 1 分钟负载为 4--14）。
迭代数不受影响；circ 和 rot45 的部分计时可能偏慢，但两种方法受到的影响相同，不会改变上面的比值结论。

\section{讨论}
\begin{itemize}
\item 文献中 AMG 的公认难题确实就在这里：不与网格对齐的各向异性破坏了经典 AMG 的强连接判据（非 M-矩阵），
插值无法沿弯曲的强方向传递信息，迭代数随 $n$ 增长（arXiv:1610.03154、arXiv:2301.13351、arXiv:2601.04112）。
\item FD-LGF 的优势与 H31--H33 一致：只要“难”的部分局限在小区域，精确局部 LU 就把它完全吸收，
剩下的背景用快速变换精确求逆，迭代数与 $n$、$\chi$ 都无关。
\item 各向异性遍布全域时，Dirichlet 边界下的旋转 9 点算子不能用正弦变换对角化（交叉项与奇延拓不交换），
FD-LGF 没有可用的精确背景，因此不适用。
\end{itemize}

\appendix
\section{迭代数随 $n$ 的完整表格}
（各表的单元为各方法在该 $(\text{系数},n)$ 下的迭代数。）
\begin{table}[H]\centering\small
\caption{迭代次数随 $n$ 的变化（$\chi=10^{2}$）}\label{tab:n2}
\begin{tabular}{lccccc}
\toprule
系数 & BAMG $\theta{=}0.25$ & BAMG $\theta{=}0.5$ & GMG & FD-LGF $w{=}\frac1{16}$ & FD-LGF $w{=}\frac1{64}$ \\
\midrule
aligned n=127 & 7 & 7 & 47 & 3 & -- \\
rot30 n=127 & 13 & 10 & 31 & 69 & -- \\
rot45 n=127 & 12 & 11 & 28 & 79 & -- \\
circ n=127 & 20 & 22 & 52 & 117 & -- \\
disk n=127 & 19 & 22 & 37 & 7 & 12 \\
band n=127 & 21 & 18 & 21 & 8 & 14 \\
aligned n=255 & 7 & 7 & 53 & 3 & -- \\
rot30 n=255 & 12 & 10 & 36 & 80 & -- \\
rot45 n=255 & 12 & 12 & 34 & 93 & -- \\
circ n=255 & 21 & 25 & 60 & 122 & -- \\
disk n=255 & 19 & 24 & 44 & 7 & 11 \\
band n=255 & 21 & 20 & 27 & 8 & 13 \\
aligned n=511 & 8 & 8 & 55 & 3 & -- \\
rot30 n=511 & 14 & 10 & 40 & 88 & -- \\
rot45 n=511 & 13 & 12 & 37 & 104 & -- \\
circ n=511 & 22 & 27 & 66 & 127 & -- \\
disk n=511 & 19 & 25 & 47 & 7 & 11 \\
band n=511 & 22 & 22 & 33 & 8 & 13 \\
aligned n=1023 & 8 & 8 & 58 & 3 & -- \\
rot30 n=1023 & 12 & 10 & 43 & 95 & -- \\
rot45 n=1023 & 13 & 13 & 40 & 112 & -- \\
circ n=1023 & 23 & 29 & 72 & 131 & -- \\
disk n=1023 & 20 & 28 & 50 & 7 & 11 \\
band n=1023 & 22 & 24 & 39 & 8 & 13 \\
\bottomrule
\end{tabular}
\end{table}

\begin{table}[H]\centering\small
\caption{迭代次数随 $n$ 的变化（$\chi=10^{4}$）}\label{tab:n4}
\begin{tabular}{lccccc}
\toprule
系数 & BAMG $\theta{=}0.25$ & BAMG $\theta{=}0.5$ & GMG & FD-LGF $w{=}\frac1{16}$ & FD-LGF $w{=}\frac1{64}$ \\
\midrule
aligned n=127 & 4 & 4 & 137 & 2 & -- \\
rot30 n=127 & 16 & 12 & 43 & 91 & -- \\
rot45 n=127 & 15 & 14 & 36 & 101 & -- \\
circ n=127 & 38 & 41 & 117 & 256 & -- \\
disk n=127 & 35 & 39 & 66 & 6 & 10 \\
band n=127 & 33 & 25 & 36 & 7 & 15 \\
aligned n=255 & 5 & 5 & 246 & 2 & -- \\
rot30 n=255 & 17 & 14 & 59 & 129 & -- \\
rot45 n=255 & 19 & 17 & 52 & 145 & -- \\
circ n=255 & 61 & 66 & 196 & 446 & -- \\
disk n=255 & 50 & 60 & 127 & 6 & 9 \\
band n=255 & 38 & 34 & 51 & 7 & 14 \\
aligned n=511 & 6 & 6 & 316 & 2 & -- \\
rot30 n=511 & 22 & 16 & 78 & 174 & -- \\
rot45 n=511 & 24 & 20 & 72 & 203 & -- \\
circ n=511 & 96 & 113 & 324 & 771 & -- \\
disk n=511 & 74 & 101 & 220 & 5 & 10 \\
band n=511 & 47 & 45 & 75 & 7 & 14 \\
aligned n=1023 & 7 & 7 & 378 & 2 & -- \\
rot30 n=1023 & 27 & 18 & 107 & 235 & -- \\
rot45 n=1023 & 33 & 25 & 101 & 278 & -- \\
circ n=1023 & 144 & 192 & 426 & 1000 & -- \\
disk n=1023 & 98 & 175 & 392 & 6 & 10 \\
band n=1023 & 56 & 62 & 101 & 7 & 14 \\
\bottomrule
\end{tabular}
\end{table}

\begin{table}[H]\centering\small
\caption{迭代次数随 $n$ 的变化（$\chi=10^{6}$）}\label{tab:n6}
\begin{tabular}{lccccc}
\toprule
系数 & BAMG $\theta{=}0.25$ & BAMG $\theta{=}0.5$ & GMG & FD-LGF $w{=}\frac1{16}$ & FD-LGF $w{=}\frac1{64}$ \\
\midrule
aligned n=127 & 2 & 2 & 101 & 2 & -- \\
rot30 n=127 & 15 & 12 & 43 & 92 & -- \\
rot45 n=127 & 15 & 14 & 36 & 102 & -- \\
circ n=127 & 40 & 43 & 120 & 263 & -- \\
disk n=127 & 38 & 42 & 68 & 4 & 7 \\
band n=127 & 34 & 25 & 44 & 5 & 11 \\
aligned n=255 & 2 & 2 & 251 & 2 & -- \\
rot30 n=255 & 17 & 14 & 59 & 130 & -- \\
rot45 n=255 & 19 & 16 & 52 & 146 & -- \\
circ n=255 & 64 & 69 & 211 & 471 & -- \\
disk n=255 & 56 & 66 & 180 & 4 & 7 \\
band n=255 & 39 & 34 & 67 & 5 & 11 \\
aligned n=511 & 2 & 2 & 503 & 2 & -- \\
rot30 n=511 & 22 & 16 & 80 & 177 & -- \\
rot45 n=511 & 25 & 18 & 74 & 206 & -- \\
circ n=511 & 108 & 128 & 384 & 874 & -- \\
disk n=511 & 85 & 113 & 301 & 4 & 7 \\
band n=511 & 48 & 46 & 108 & 6 & 10 \\
aligned n=1023 & 3 & 3 & 1000 & 2 & -- \\
rot30 n=1023 & 28 & 19 & 111 & 244 & -- \\
rot45 n=1023 & 33 & 21 & 104 & 287 & -- \\
circ n=1023 & 191 & 241 & 718 & 1000 & -- \\
disk n=1023 & 136 & 218 & 732 & 4 & 7 \\
band n=1023 & 59 & 65 & 167 & 6 & 11 \\
\bottomrule
\end{tabular}
\end{table}
\end{document}
